\documentclass[11pt]{article}

\usepackage[margin=1in]{geometry}
\usepackage{amsmath,amssymb}
\usepackage{graphicx}
\usepackage{booktabs}
\usepackage{array}
\usepackage{siunitx}
\usepackage{hyperref}
\usepackage{xcolor}
\usepackage{enumitem}
\usepackage{longtable}
\usepackage{float}

\hypersetup{
    colorlinks=true,
    linkcolor=blue!60!black,
    citecolor=blue!60!black,
    urlcolor=blue!60!black
}

\sisetup{detect-all=true, per-mode=symbol}
\setlist[itemize]{topsep=2pt,itemsep=2pt,parsep=0pt}
\setlist[enumerate]{topsep=2pt,itemsep=2pt,parsep=0pt}

\title{\textbf{Taylor Cone Experimental Validation:\\Detailed Apparatus and Procedure}}
\author{Ethan Yang, Elliot Dong, Curtis Lau\\
Independent Schools Foundation Academy\\
Hong Kong SAR}
\date{\today}

\begin{document}

\maketitle

\tableofcontents
\newpage

\section{Introduction}

This document provides complete experimental apparatus specifications, detailed safety protocols, and step-by-step procedures for Taylor cone onset validation experiments. It is designed as a standalone reference for experimental implementation of the theoretical and computational work described in the main ISEF paper.

\subsection{Document Purpose}

\begin{itemize}
\item Complete equipment list with budget estimates and suppliers
\item Detailed safety protocols addressing all hazards
\item Step-by-step setup and operational procedures
\item Troubleshooting guide and emergency procedures
\item Required for institutional safety review and approval
\end{itemize}

\subsection{Scope}

This procedure covers:
\begin{itemize}
\item High-voltage electrospray apparatus assembly
\item Camera-based Taylor cone imaging
\item Voltage ramp experiments for onset detection
\item Image processing and data analysis workflow
\end{itemize}

\section{Complete Apparatus List with Budget}

\subsection{High Voltage Power Supply}

\paragraph{Primary Option: Current-Limited Lab Supply}

\begin{itemize}
\item \textbf{Item:} Glassman High Voltage PS/EH Series (0--10~kV, 1--2~mA)
\item \textbf{Specifications:}
  \begin{itemize}
  \item Output: 0--10,000~V DC
  \item Current limit: 1--2~mA (critical for safety)
  \item Regulation: $\pm0.01\%$
  \item Digital display and remote control
  \end{itemize}
\item \textbf{Cost:} \$1,500--2,500 USD
\item \textbf{Supplier:} Glassman High Voltage, Newark Electronics, Mouser
\item \textbf{Safety features:} Hardware current limiting, emergency shutdown, discharge circuit
\end{itemize}

\paragraph{Alternative Option: More Affordable}

\begin{itemize}
\item \textbf{Item:} Spellman SL Series (0--5~kV, 1~mA)
\item \textbf{Cost:} \$800--1,200 USD
\item \textbf{Supplier:} Spellman HV Electronics, eBay (used units)
\end{itemize}

\paragraph{Budget Option: Regulated Flyback Module}

\begin{itemize}
\item \textbf{Item:} XP Power CA/CB Series with external current limiter
\item \textbf{Cost:} \$200--400 USD
\item \textbf{Supplier:} DigiKey, Mouser
\item \textbf{Note:} Requires additional enclosure and safety circuitry
\end{itemize}

\textbf{Recommended:} Glassman or Spellman for lab safety compliance.

\subsection{Syringe Pump}

\paragraph{Primary Option: Precision Syringe Pump}

\begin{itemize}
\item \textbf{Item:} Harvard Apparatus PHD 2000 or PHD ULTRA
\item \textbf{Specifications:}
  \begin{itemize}
  \item Flow rate range: 0.001~$\mu$L/h to 220~mL/h
  \item Accuracy: $\pm0.5\%$
  \item Compatible with 0.5--60~mL syringes
  \end{itemize}
\item \textbf{Cost:} \$2,000--3,500 USD
\item \textbf{Supplier:} Harvard Apparatus, VWR
\end{itemize}

\paragraph{Alternative Option: More Affordable}

\begin{itemize}
\item \textbf{Item:} New Era NE-300 Single Syringe Pump
\item \textbf{Specifications:}
  \begin{itemize}
  \item Flow rate: 0.73~$\mu$L/h to 2120~mL/h
  \item Accuracy: $\pm1\%$
  \end{itemize}
\item \textbf{Cost:} \$800--1,200 USD
\item \textbf{Supplier:} New Era Pump Systems (syringepump.com)
\end{itemize}

\paragraph{Budget Option: DIY Stepper Motor Pump}

\begin{itemize}
\item \textbf{Item:} Stepper motor + lead screw + Arduino control
\item \textbf{Cost:} \$150--300 USD
\item \textbf{Supplier:} Parts from Amazon, Adafruit
\item \textbf{Note:} Requires calibration and flow verification
\end{itemize}

\textbf{Recommended:} New Era NE-300 (good balance of cost and precision).

\subsection{Camera and Imaging System}

\paragraph{Primary Option: High-Speed CMOS Camera}

\begin{itemize}
\item \textbf{Item:} FLIR Blackfly S USB3 (BFS-U3-32S4M-C)
\item \textbf{Specifications:}
  \begin{itemize}
  \item Resolution: 2048$\times$1536 pixels
  \item Frame rate: 118~fps at full resolution
  \item Sensor: Sony IMX252 (3.45~$\mu$m pixel)
  \item Interface: USB 3.0
  \end{itemize}
\item \textbf{Cost:} \$600--900 USD
\item \textbf{Supplier:} FLIR, Edmund Optics
\end{itemize}

\paragraph{Alternative Option: Mid-Range USB Camera}

\begin{itemize}
\item \textbf{Item:} ELP USB Camera (5MP, 2592$\times$1944)
\item \textbf{Specifications:}
  \begin{itemize}
  \item Resolution: 5MP
  \item Frame rate: 15--30~fps
  \item Manual focus
  \end{itemize}
\item \textbf{Cost:} \$80--150 USD
\item \textbf{Supplier:} Amazon, AliExpress
\end{itemize}

\paragraph{Budget Option: Webcam with Macro Lens}

\begin{itemize}
\item \textbf{Item:} Logitech C920 + macro lens adapter
\item \textbf{Cost:} \$100--150 USD
\item \textbf{Supplier:} Amazon, Best Buy
\item \textbf{Note:} Limited resolution but sufficient for proof-of-concept
\end{itemize}

\textbf{Recommended:} FLIR Blackfly S (research-grade) or ELP camera (budget).

\subsection{Budget Summary}

\begin{table}[H]
\centering
\caption{Research-Grade Setup (\$7,000--10,000 USD)}
\small
\begin{tabular}{lr}
\toprule
Component & Cost (USD) \\
\midrule
Glassman HV Power Supply & \$2,000 \\
Harvard Apparatus Syringe Pump & \$2,500 \\
FLIR Blackfly Camera & \$800 \\
Telecentric Lens & \$1,500 \\
LED Backlight & \$150 \\
Optical Breadboard Setup & \$500 \\
Metal Enclosure + Safety & \$500 \\
Needles, Electrodes, Hardware & \$300 \\
Chemicals and Consumables & \$200 \\
PPE and Safety Equipment & \$200 \\
Calibration Tools & \$150 \\
Miscellaneous & \$200 \\
\midrule
\textbf{Total} & \textbf{\$9,000} \\
\bottomrule
\end{tabular}
\end{table}

\begin{table}[H]
\centering
\caption{Mid-Range Setup (\$3,000--4,000 USD)}
\small
\begin{tabular}{lr}
\toprule
Component & Cost (USD) \\
\midrule
Spellman HV Power Supply & \$1,000 \\
New Era NE-300 Syringe Pump & \$1,000 \\
ELP USB Camera & \$120 \\
C-Mount Macro Lens + Extensions & \$80 \\
Huion LED Light Pad & \$50 \\
Optical Breadboard Setup & \$400 \\
Metal Enclosure + Safety & \$400 \\
Needles, Electrodes, Hardware & \$250 \\
Chemicals and Consumables & \$150 \\
PPE and Safety Equipment & \$150 \\
Calibration Tools & \$100 \\
Miscellaneous & \$150 \\
\midrule
\textbf{Total} & \textbf{\$3,850} \\
\bottomrule
\end{tabular}
\end{table}

\begin{table}[H]
\centering
\caption{Budget Setup (\$1,200--1,800 USD)}
\small
\begin{tabular}{lr}
\toprule
Component & Cost (USD) \\
\midrule
XP Power HV Module + Limiter & \$300 \\
DIY Arduino Syringe Pump & \$200 \\
Logitech C920 + Macro Adapter & \$130 \\
DIY LED Panel & \$35 \\
DIY Aluminum Frame & \$80 \\
Acrylic Box + Wire Mesh & \$180 \\
Needles, Electrodes, Hardware & \$150 \\
Chemicals and Consumables & \$120 \\
PPE and Safety Equipment & \$120 \\
Digital Calipers & \$60 \\
Miscellaneous & \$100 \\
\midrule
\textbf{Total} & \textbf{\$1,475} \\
\bottomrule
\end{tabular}
\end{table}

\section{Safety Protocols and Risk Mitigation}

\subsection{Primary Hazards}

\begin{enumerate}
\item \textbf{Electric shock} (potentially lethal)
\item \textbf{High voltage arc/spark} (fire risk)
\item \textbf{Chemical exposure} (ethanol, solvents)
\item \textbf{Sharp needle} (puncture injury)
\end{enumerate}

\subsection{High Voltage Safety Protocol}

\subsubsection{Before ANY Experiment}

\textbf{Required checks:}

\begin{itemize}
\item[$\square$] All participants have read and signed safety acknowledgment
\item[$\square$] Supervisor present and briefed on emergency procedures
\item[$\square$] Emergency contact numbers posted on wall
\item[$\square$] Fire extinguisher (Class C) within 3 meters
\item[$\square$] First aid kit accessible
\item[$\square$] Workspace clear of flammable materials
\item[$\square$] All personnel wearing safety glasses
\item[$\square$] Electrical insulating gloves available
\item[$\square$] Door interlock tested and functional
\item[$\square$] Emergency stop button tested
\item[$\square$] All ground connections verified ($<1\Omega$ to earth)
\item[$\square$] Current limit verified at 1~mA
\item[$\square$] Only ONE person operates HV controls
\item[$\square$] All observers maintain 2~m distance from enclosure
\end{itemize}

\subsubsection{During Experiment}

\textbf{Operating rules:}

\begin{itemize}
\item Door MUST be closed before applying HV
\item NO adjustments while HV is on
\item If adjustment needed: TURN OFF HV $\rightarrow$ wait 60s $\rightarrow$ open door $\rightarrow$ adjust
\item Only HV operator touches controls
\item All observers stay behind marked line (2~m from enclosure)
\item Monitor for unusual sounds (buzzing, crackling) $\rightarrow$ shut down immediately
\item Monitor for unusual smells (ozone, burning) $\rightarrow$ shut down immediately
\end{itemize}

\subsubsection{After Experiment}

\textbf{Shutdown procedure:}

\begin{itemize}
\item Turn voltage knob to zero (fully CCW)
\item Turn off HV power supply main switch
\item Wait minimum 60 seconds (capacitor discharge time)
\item Open door (safe to enter)
\item Use insulated stick to short HV terminal to ground (extra precaution)
\item Only then touch equipment
\item Disconnect HV cable from needle
\item Cover exposed HV terminals with insulating caps
\end{itemize}

\subsubsection{Emergency: Someone is Being Shocked}

\begin{center}
\textcolor{red}{\textbf{DO NOT TOUCH THE PERSON}}
\end{center}

\begin{enumerate}
\item \textbf{Press emergency stop button} (breaks HV circuit)
\item \textbf{If E-stop doesn't work:} Unplug HV power supply
\item \textbf{Wait 5 seconds} (let capacitors discharge)
\item \textbf{Only then:} Check victim, provide first aid, call emergency services
\end{enumerate}

\textbf{Hong Kong Emergency:} 999 (Fire, Ambulance, Police)

\subsection{Chemical Safety Protocol}

\subsubsection{Ethanol Handling}

\begin{itemize}
\item \textbf{Hazard:} Flammable, eye irritant
\item \textbf{PPE required:} Safety glasses, nitrile gloves, lab coat
\item \textbf{Ventilation:} Use in well-ventilated area or fume hood
\item \textbf{Storage:} In flammable-safe cabinet, away from HV equipment
\item \textbf{Spill cleanup:} Absorb with paper towels, dispose in hazardous waste
\item \textbf{Disposal:} Follow local hazardous waste regulations (do NOT pour down drain)
\end{itemize}

\subsubsection{First Aid}

\begin{itemize}
\item \textbf{Eye contact:} Rinse immediately with water for 15 minutes, seek medical attention
\item \textbf{Skin contact:} Wash with soap and water
\item \textbf{Ingestion:} Do not induce vomiting, seek medical attention immediately
\item \textbf{Inhalation:} Move to fresh air
\end{itemize}

\subsection{Current Limit Verification}

\textbf{CRITICAL SAFETY STEP:}

\begin{enumerate}
\item \textbf{Set current limit to 1~mA:}
  \begin{itemize}
  \item Adjust current limit pot/knob to minimum
  \item Use multimeter in current measurement mode if possible
  \end{itemize}

\item \textbf{Why 1~mA limit is critical:}
  \begin{itemize}
  \item Current through human body:
    \begin{itemize}
    \item $>1$~mA: Barely perceptible
    \item $>10$~mA: Painful, cannot let go
    \item $>100$~mA: Potentially lethal (ventricular fibrillation)
    \end{itemize}
  \item 1~mA limit provides safety margin
  \item Still sufficient for Taylor cone formation (nA--$\mu$A typical)
  \end{itemize}
\end{enumerate}

\section{Detailed Setup Procedure}

\subsection{Phase 1: Pre-Assembly Preparation (1--2 days)}

\subsubsection{Workspace Preparation}

\paragraph{Select appropriate workspace:}

\begin{itemize}
\item Lab bench or sturdy table (minimum 1.2~m $\times$ 0.8~m)
\item Near grounded electrical outlet
\item Good lighting
\item Fire extinguisher accessible (Class C for electrical fires)
\item First aid kit nearby
\item Clear of flammable materials
\end{itemize}

\paragraph{Verify electrical safety:}

\begin{itemize}
\item Confirm outlet is properly grounded (use outlet tester)
\item Check circuit breaker rating (minimum 10~A recommended)
\item Test ground continuity (should read $<1\Omega$ to earth)
\end{itemize}

\subsection{Phase 2: Mechanical Assembly (2--3 hours)}

\subsubsection{Enclosure Setup}

\begin{enumerate}
\item \textbf{Position enclosure on bench:}
  \begin{itemize}
  \item Leave 30~cm clearance on all sides for access
  \item Ensure door can open fully
  \item Mark locations for cable entry
  \end{itemize}

\item \textbf{Ground the enclosure:}
  \begin{itemize}
  \item Attach ground wire to metal frame
  \item Run to building ground or ground rod
  \item Verify continuity: $<1\Omega$ to earth ground
  \item Label: ``SAFETY GROUND - DO NOT REMOVE''
  \end{itemize}

\item \textbf{Install interlock switch:}
  \begin{itemize}
  \item Mount magnetic switch on door frame
  \item Wire normally-open contact to HV supply enable line
  \item Test: Opening door should cut HV output immediately
  \item Label switch: ``SAFETY INTERLOCK''
  \end{itemize}

\item \textbf{Install emergency stop button:}
  \begin{itemize}
  \item Mount red mushroom button outside enclosure
  \item Wire in series with HV enable
  \item Test functionality before proceeding
  \item Label: ``EMERGENCY STOP''
  \end{itemize}
\end{enumerate}

\subsubsection{Optical Breadboard Assembly}

\begin{enumerate}
\item \textbf{Mount breadboard inside enclosure:}
  \begin{itemize}
  \item Use rubber feet or vibration dampers
  \item Ensure level (use spirit level)
  \item Leave space for HV supply underneath or outside
  \end{itemize}

\item \textbf{Install optical posts:}
  \begin{itemize}
  \item Post 1: Camera mount (rear, centered)
  \item Post 2: Needle/syringe holder (front-center)
  \item Post 3--4: Electrode mount (front, vertically adjustable)
  \item Tighten securely but avoid over-torquing
  \end{itemize}
\end{enumerate}

\subsection{Phase 3: Electrical Connections (1--2 hours)}

\begin{center}
\textcolor{red}{\textbf{CRITICAL: All electrical work must be done with HV power supply UNPLUGGED}}
\end{center}

\subsubsection{HV Power Supply Placement}

\begin{enumerate}
\item \textbf{Position HV supply:}
  \begin{itemize}
  \item \textbf{Outside} enclosure if possible (easier access, better cooling)
  \item On rubber mat (electrical insulation)
  \item Away from liquids
  \item Ventilation unobstructed
  \end{itemize}

\item \textbf{Verify all controls are at minimum/off:}
  \begin{itemize}
  \item Voltage knob: Minimum (fully CCW)
  \item Current limit: Minimum
  \item Power switch: OFF
  \item Main power: UNPLUGGED
  \end{itemize}
\end{enumerate}

\subsubsection{HV Output Connection}

\begin{enumerate}
\item \textbf{Select appropriate HV cable:}
  \begin{itemize}
  \item Use HV-rated silicone insulated wire (20--30~kV rating)
  \item Length: Minimize excess (2--3~m maximum)
  \item Color: RED for high voltage
  \item Do NOT use standard hookup wire
  \end{itemize}

\item \textbf{Connect HV output to needle:}
  \begin{itemize}
  \item Route HV cable through enclosure grommet
  \item Connect to needle holder (must be conductive)
  \item Use HV-rated alligator clip or crimp lug
  \item Ensure connection is mechanically secure
  \item Wrap connection with HV-rated electrical tape
  \item Keep cable away from grounded surfaces ($>5$~cm clearance)
  \end{itemize}

\item \textbf{Label HV cable:}
  \begin{itemize}
  \item Attach warning labels every 30~cm
  \item ``HIGH VOLTAGE - DO NOT TOUCH''
  \item Include maximum voltage rating
  \end{itemize}
\end{enumerate}

\subsubsection{Ground Connections}

\begin{enumerate}
\item \textbf{Ground the extractor electrode:}
  \begin{itemize}
  \item Connect thick ground wire (12 AWG minimum)
  \item Run to common ground point
  \item Verify continuity to building ground
  \end{itemize}

\item \textbf{Create single-point ground:}
  \begin{itemize}
  \item All grounds should connect at ONE point
  \item This prevents ground loops
  \item Use ground bus bar or large terminal block
  \item Final connection to building ground or ground rod
  \end{itemize}
\end{enumerate}

\section{Step-by-Step Experimental Protocol}

\subsection{Pre-Experiment Preparation (30 minutes)}

\subsubsection{Day Before Experiment}

\begin{enumerate}
\item \textbf{Check equipment:}
  \begin{itemize}
  \item Verify all components are functional
  \item Check HV cable insulation (no cracks or damage)
  \item Test camera and image capture
  \item Verify sufficient liquid supply
  \end{itemize}

\item \textbf{Prepare test liquid:}
  \begin{itemize}
  \item For ethanol + NaCl:
    \begin{itemize}
    \item Measure 100~mL ethanol in graduated cylinder
    \item Dissolve 0.01--0.1~g NaCl (adjust for desired conductivity)
    \item Stir until fully dissolved
    \item Store in sealed bottle, label clearly
    \item Record: date, concentration, operator name
    \end{itemize}
  \end{itemize}

\item \textbf{Clean equipment:}
  \begin{itemize}
  \item Wipe down electrode with ethanol (remove residue)
  \item Inspect needle for contamination
  \item Clean optical surfaces (lens, backlight diffuser)
  \end{itemize}
\end{enumerate}

\subsection{Voltage Ramp Protocol (30--60 minutes)}

\begin{enumerate}
\item \textbf{Initial state:}
  \begin{itemize}
  \item Voltage: 0~V
  \item Pump: OFF
  \item Camera: Recording
  \item Note starting timestamp
  \end{itemize}

\item \textbf{Start liquid flow:}
  \begin{itemize}
  \item Set pump to 0.5--2~$\mu$L/min
  \item Observe meniscus forming at needle tip
  \item Wait until steady hemispherical droplet appears ($\sim$1--2~min)
  \item If droplet falls: reduce flow rate
  \item Record stable flow rate in metadata
  \end{itemize}

\item \textbf{Close enclosure and enable HV:}
  \begin{itemize}
  \item Close door (verify interlock engages)
  \item Turn on HV power supply main switch
  \item Verify enable LED is lit
  \item Voltage should still be at zero
  \end{itemize}

\item \textbf{Begin voltage sweep:}
  \begin{itemize}
  \item \textbf{Step size:} 100--200~V per step
  \item \textbf{Dwell time:} 20--30 seconds per step
  \item \textbf{Procedure for each step:}
    \begin{enumerate}
    \item Increase voltage by one step
    \item Read voltage display, announce/record value
    \item Observe meniscus through camera feed
    \item Classify state:
      \begin{itemize}
      \item ``Hemispherical'' - no deformation
      \item ``Ellipsoidal'' - slight elongation
      \item ``Conical-unstable'' - cone forms but oscillates
      \item ``Conical-stable'' - steady cone, no spray
      \item ``Spray'' - visible jet/droplet emission
      \end{itemize}
    \item Make verbal annotation (recorded by camera audio)
    \item Wait for full dwell time before next step
    \end{enumerate}
  \end{itemize}

\item \textbf{Identify onset condition:}
  \begin{itemize}
  \item \textbf{Onset voltage ($V_0$):} First voltage where stable cone appears
  \item Hold at this voltage for 60 seconds
  \item Verify stability (no jetting, no collapse)
  \item Record voltage and timestamp
  \end{itemize}

\item \textbf{STOP immediately if:}
  \begin{itemize}
  \item Sparking/arcing occurs $\rightarrow$ shut down, investigate
  \item Unusual odor $\rightarrow$ shut down, ventilate
  \item Current meter shows $>1$~mA $\rightarrow$ shut down, check connections
  \end{itemize}
\end{enumerate}

\subsection{Shutdown Procedure (10 minutes)}

\begin{enumerate}
\item \textbf{Power down:}
  \begin{itemize}
  \item Turn voltage to zero
  \item Turn off HV power supply
  \item Wait 60 seconds
  \item Stop syringe pump
  \item Stop video recording
  \end{itemize}

\item \textbf{Safe equipment access:}
  \begin{itemize}
  \item Open enclosure door
  \item Short HV terminal to ground (insulated stick)
  \item Disconnect HV cable from needle
  \item Remove test liquid syringe
  \end{itemize}

\item \textbf{Clean up:}
  \begin{itemize}
  \item Wipe any liquid residue from electrode
  \item Dispose of used needle in sharps container
  \item Cap and label used liquid (if reusing)
  \item Store liquids in flammable cabinet
  \end{itemize}

\item \textbf{Data backup:}
  \begin{itemize}
  \item Copy video files to backup drive
  \item Verify file integrity (can play back)
  \item Fill in metadata template
  \item Save metadata as JSON file with same filename as video
  \end{itemize}
\end{enumerate}

\section{Troubleshooting Guide}

\subsection{Problem 1: No meniscus forms at needle tip}

\textbf{Possible causes:}
\begin{itemize}
\item Flow rate too low
\item Needle clogged
\item Liquid wetting properties poor
\end{itemize}

\textbf{Solutions:}
\begin{itemize}
\item Increase flow rate gradually (0.5 $\rightarrow$ 1 $\rightarrow$ 2~$\mu$L/min)
\item Check for blockage: disconnect needle, run pump, verify flow
\item Replace needle if clogged
\item Check liquid surface tension (should be $>20$~mN/m for stable meniscus)
\end{itemize}

\subsection{Problem 2: No cone formation even at high voltage}

\textbf{Possible causes:}
\begin{itemize}
\item Liquid not conductive enough
\item Electrode spacing too large
\item Voltage too low
\item Poor electrical connection
\end{itemize}

\textbf{Solutions:}
\begin{itemize}
\item Increase dopant concentration (add more NaCl)
\item Measure liquid conductivity (should be $>10$~$\mu$S/cm)
\item Reduce electrode spacing to 10--15~mm
\item Check HV connection to needle (verify continuity)
\item Increase voltage gradually (literature onset: 2--5~kV typical)
\end{itemize}

\subsection{Problem 3: Sparking/arcing between needle and electrode}

\textbf{Possible causes:}
\begin{itemize}
\item Voltage too high for given spacing
\item Sharp edges on electrode
\item Contamination on surfaces
\end{itemize}

\textbf{Solutions:}
\begin{itemize}
\item \textbf{Immediately reduce voltage to zero}
\item Increase electrode spacing
\item Polish electrode edges (remove burrs)
\item Clean surfaces with ethanol
\item Check for conductive debris (liquid droplets on surfaces)
\end{itemize}

\section{Emergency Procedures}

\subsection{Emergency 1: Electrical Shock}

\textbf{If someone is being shocked:}

\begin{enumerate}
\item \textbf{DO NOT TOUCH THE PERSON}
\item Press emergency stop button or unplug HV supply
\item Wait 5 seconds (let capacitors discharge)
\item Check victim:
  \begin{itemize}
  \item Conscious? $\rightarrow$ Reassure, check for burns, seek medical attention
  \item Unconscious? $\rightarrow$ Call 999 immediately, begin CPR if trained
  \end{itemize}
\item Do not move victim unless necessary for safety
\end{enumerate}

\subsection{Emergency 2: Fire}

\textbf{Small fire ($<$ 30~cm):}

\begin{enumerate}
\item Press emergency stop (cut power)
\item Use fire extinguisher (Class C)
\item Aim at base of fire, sweep side to side
\item If extinguished: ventilate area, document incident
\item If spreads: proceed to large fire procedure
\end{enumerate}

\textbf{Large fire or spreading rapidly:}

\begin{enumerate}
\item Evacuate immediately (do not stop for belongings)
\item Close door to lab
\item Pull fire alarm
\item Call 999 from safe location
\item Meet at designated assembly point
\item Do NOT re-enter building
\end{enumerate}

\subsection{Emergency 3: Chemical Spill}

\textbf{Small spill ($<$ 100~mL):}

\begin{enumerate}
\item Alert others in area
\item Put on gloves and safety glasses
\item Absorb with paper towels or spill kit absorbent
\item Dispose in chemical waste container
\item Wipe area with water
\item Ventilate area
\end{enumerate}

\textbf{Spill on person:}

\begin{itemize}
\item Remove contaminated clothing immediately
\item Rinse affected area with water for 15 minutes
\item Seek medical attention if irritation persists
\end{itemize}

\section{Pre-Experiment Safety Checklist}

\textbf{Print and complete before EVERY experiment session:}

\vspace{0.5cm}

\noindent\textbf{Date:} \underline{\hspace{3cm}}  
\textbf{Operators:} \underline{\hspace{10cm}}  
\textbf{Supervisor:} \underline{\hspace{4cm}}  
\textbf{Start time:} \underline{\hspace{2cm}}

\subsection{Equipment Checks}

\begin{itemize}
\item[$\square$] HV power supply: voltage at zero, current limit at 1~mA
\item[$\square$] Interlock switch: tested and functional (open door $\rightarrow$ HV disables)
\item[$\square$] Emergency stop: tested and functional (press button $\rightarrow$ HV disables)
\item[$\square$] Ground connections: verified $<1\Omega$ to earth
\item[$\square$] HV cable: inspected for damage (no cracks, no exposed conductor)
\item[$\square$] Enclosure: door closes properly, no gaps
\item[$\square$] Camera: recording function tested
\item[$\square$] Backlight: functional, no flickering
\item[$\square$] Syringe pump: smooth operation, no leaks
\item[$\square$] Liquid supply: sufficient for planned experiments, properly labeled
\end{itemize}

\subsection{Safety Equipment}

\begin{itemize}
\item[$\square$] Fire extinguisher: present, charged, within 3 meters
\item[$\square$] First aid kit: present and accessible
\item[$\square$] Safety glasses: all participants wearing
\item[$\square$] Electrical insulating gloves: available at workstation
\item[$\square$] Nitrile gloves: available
\item[$\square$] Lab coats: all participants wearing
\end{itemize}

\subsection{Workspace}

\begin{itemize}
\item[$\square$] Area cleared of flammable materials ($>1$~m from enclosure)
\item[$\square$] Emergency contact numbers posted on wall
\item[$\square$] Adequate ventilation (windows open or fume hood on)
\item[$\square$] Clear path to exit (no obstructions)
\item[$\square$] Observers maintain 2~m distance (marked line on floor)
\end{itemize}

\subsection{Personnel}

\begin{itemize}
\item[$\square$] All participants have read safety protocols
\item[$\square$] All participants have signed safety acknowledgment
\item[$\square$] Roles assigned (HV operator: \underline{\hspace{3cm}}, recorder: \underline{\hspace{3cm}})
\item[$\square$] Everyone knows location of E-stop and fire extinguisher
\item[$\square$] Everyone knows emergency phone number: 999
\end{itemize}

\subsection{Final Checks}

\begin{itemize}
\item[$\square$] Supervisor present and briefed
\item[$\square$] Mobile phone present (for emergency calls)
\item[$\square$] No food or drinks in work area
\item[$\square$] No loose clothing or jewelry near moving parts
\end{itemize}

\vspace{1cm}

\noindent\textbf{Supervisor signature:} \underline{\hspace{5cm}}  
\textbf{Approval to proceed:} YES / NO

\section{Appendix: Quick Reference}

\subsection{Onset Voltage Estimation}

Use this to predict expected onset voltage for your geometry:

\begin{equation}
V_0 \approx \sqrt{\frac{4\gamma L}{\varepsilon_0 \ln(4L/r_{\text{needle}})}}
\end{equation}

Where:
\begin{itemize}
\item $\gamma$ = surface tension (N/m)
\item $L$ = electrode spacing (m)
\item $\varepsilon_0 = 8.854 \times 10^{-12}$~F/m
\item $r_{\text{needle}}$ = needle outer radius (m)
\end{itemize}

\textbf{Example calculation (ethanol, 21-gauge needle, 20~mm spacing):}

\begin{itemize}
\item $\gamma = 0.022$~N/m (ethanol)
\item $L = 0.020$~m
\item $r_{\text{needle}} = 0.0004$~m (21-gauge OD = 0.8~mm)
\item $\ln(4 \times 0.020/0.0004) = \ln(200) \approx 5.3$
\end{itemize}

\begin{align}
V_0 &\approx \sqrt{\frac{4 \times 0.022 \times 0.020}{8.854 \times 10^{-12} \times 5.3}} \\
&\approx \sqrt{\frac{0.00176}{4.69 \times 10^{-11}}} \\
&\approx \sqrt{3.75 \times 10^7} \\
&\approx 6,100~\text{V}
\end{align}

\textbf{Expect onset around 6~kV for this geometry.}

\subsection{Hong Kong Supplier Contacts}

\paragraph{Fisher Scientific Hong Kong}
\begin{itemize}
\item Phone: +852 2407 2600
\item Website: fishersci.com.hk
\item Products: Chemicals, lab equipment, consumables
\end{itemize}

\paragraph{VWR International (Asia Pacific)}
\begin{itemize}
\item Phone: +852 2922 2082
\item Website: hk.vwr.com
\item Products: Lab equipment, chemicals
\end{itemize}

\paragraph{RS Components Hong Kong}
\begin{itemize}
\item Phone: +852 2610 3888
\item Website: hk.rs-online.com
\item Products: Test equipment, tools, components
\end{itemize}

\paragraph{Sham Shui Po Electronics Market}
\begin{itemize}
\item Location: Apliu Street (鴨寮街), Sham Shui Po
\item MTR: Sham Shui Po Station Exit C2
\item Hours: Daily 10:00--20:00
\item Best for: Cables, connectors, tools, basic electronics
\end{itemize}

\section{Document Version Control}

\textbf{Version:} 1.0  
\textbf{Date:} \today  
\textbf{Authors:} Ethan Yang, Elliot Dong, Curtis Lau  
\textbf{Institution:} Independent Schools Foundation Academy  
\textbf{Reviewed by:} \underline{\hspace{5cm}}  
\textbf{Approval date:} \underline{\hspace{3cm}}

\end{document}
